% examples/university-example.tex
%
% Document universitaire malgache d'exemple pour babel-malagasy.
% Sert de test d'intégration: page de titre, table des matières,
% chapitres, sections, figure, tableau, bibliographie, document
% bilingue malgache/français.
%
% STATUT: ce document a été réellement compilé avec pdflatex
% pendant la préparation du projet (voir le rapport). Il ne
% requiert que ce que babel-malagasy fournit réellement:
% \babelprovide[import]{malagasy} + \setlocalecaption, chargés
% via \usepackage{babel-malagasy}.
%
\documentclass[12pt]{report}
\usepackage[T1]{fontenc}
\usepackage[utf8]{inputenc}
\usepackage{babel-malagasy}
\babelprovide[import]{french} % see docs/linguistic-policy.md:
  % this environment's TeX Live install has no "french" babel
  % package option, but the ini locale loads fine directly.
\usepackage{graphicx}
\usepackage{booktabs}

\title{Ny Fikarohana momba ny Fampiasana ny Teknolojia \\
  ao amin'ny Fampianarana Ambaratonga Ambony \\
  \large (Recherche sur l'usage des technologies \\
  dans l'enseignement supérieur)}
\author{Anarana Mpikaroka}
\date{\today}

\begin{document}
\selectlanguage{malagasy}

\maketitle

\tableofcontents
\listoffigures
\listoftables

\chapter{Fampidirana}

Ity boky ity dia mikasika ny \emph{fampiasana ny teknolojia} ao
amin'ny \emph{fampianarana ambaratonga ambony} eto Madagasikara.
Ny \emph{fikarohana} natao dia nifantoka tamin'ny \emph{oniversite}
telo, ary ny \emph{mpianatra} sy ny \emph{mpampianatra} no
nandray anjara tamin'izany.

Andro anio: \today.

\section{Tanjona}

Ny tanjon'ny \emph{fanatanterahana} ity \emph{fikarohana} ity dia
ny hamantarana ny \emph{fampandrosoana} azo atao amin'ny
\emph{fifandraisana} eo amin'ny \emph{sekoly} sy ny \emph{tany}
manodidina azy, indrindra ao amin'ny sehatry ny
\emph{informatika}, ny \emph{matematika}, ary ny
\emph{statistika}.

\chapter{Ny Fomba Fikarohana}

Ny \emph{fanabeazana} momba ny \emph{teknolojia} dia mila
\emph{fikarohana} lalina. Nampiasaina teo amin'ity asa ity ny
\emph{programa} maromaro mba hamakafakana ny \emph{boky} sy ny
\emph{sary} nangonina tany amin'ny \emph{trano} fampianarana
isan-karazany. Ny \emph{rano} sy ny \emph{tany}, ho ohatra ihany,
dia tsy nampiasaina afa-tsy ho fanoharana momba ny fototeny
malagasy tsotra ampiasaina amin'ny fitsapana ny fizarana teny.

\begin{figure}[h]
  \centering
  \fbox{\rule{0pt}{3cm}\rule{6cm}{0pt}}
  \caption{Ohatra sary (placeholder) mampiseho ny \emph{olona}
    mianatra \emph{algoritma} tsotra.}
\end{figure}

\begin{table}[h]
  \centering
  \begin{tabular}{lcc}
    \toprule
    Oniversite & Mpianatra & Mpampianatra \\
    \midrule
    Oniversite A & 120 & 8 \\
    Oniversite B & 95  & 6 \\
    Oniversite C & 140 & 10 \\
    \bottomrule
  \end{tabular}
  \caption{Isan'ny mpianatra sy mpampianatra araka ny oniversite.}
\end{table}

\chapter{Fehin-kevitra}

Ny \emph{fikarohana} dia nampiseho fa mbola misy ny
\emph{fampandrosoana} tokony hatao amin'ny \emph{fanatanterahana}
ny \emph{teknolojia} ao amin'ny \emph{fampianarana}. Toy izao no
tokony ho fintinina:

\begin{abstract}
Ity \emph{famintinana} ity dia manamarina fa ny
\emph{fampiasana} ny \emph{teknolojia} ao amin'ny
\emph{fampianarana} ambaratonga ambony eto Madagasikara dia
mbola mitaky \emph{fikarohana} bebe kokoa.
\end{abstract}

\appendix
\chapter{Fanampiny statistika}

Jereo ny \tablename~2.1 etsy ambony.

\selectlanguage{french}
\chapter*{Résumé en français}
\addcontentsline{toc}{chapter}{Résumé en français}

Ce chapitre, en français, illustre le passage d'une langue à
l'autre au sein d'un même document, via
\texttt{\textbackslash selectlanguage}. \foreignlanguage{malagasy}{Ary
ity fehezanteny ity dia amin'ny teny malagasy}, insérée localement
avec \texttt{\textbackslash foreignlanguage}.

\selectlanguage{malagasy}

\begin{thebibliography}{9}
\bibitem{ex1} Anarana Mpanoratra, \emph{Lohateny boky},
  Antananarivo, 2024.
\end{thebibliography}

\end{document}
